% ID: SEL-D0053-OA-E03
% ESTADO: cotejada. Segunda pasada visual y borrador compilado revisados.
% Fuente: Selectividad/sel_2013_matematicas_aplicadas/septiembre/titular/EXAMEN 2-ANDALUCÍA.pdf
% SHA-256: f8633bbd91d33bee02c2cc582c33b7ffeb3643f312b810ad0bb8c5d154d0721d
% Página: 1; opción A; ejercicio original 3.

Se cree que hay una vuelta hacia estilos de baile más populares, por lo que se realiza una encuesta a estudiantes de bachillerato, resultando que al 40\% les gusta la salsa, al 30\% les gusta el merengue y al 10\% les gusta tanto la salsa como el merengue.
\begin{enumerate}[label=\alph*)]
\item \textbf{(0.75 puntos)} \textquestiondown{}Cuál es la probabilidad de que a un estudiante le guste el merengue si le gusta la salsa?
\item \textbf{(0.75 puntos)} \textquestiondown{}Y la de que a un estudiante le guste el merengue si no le gusta la salsa?
\item \textbf{(1 punto)} \textquestiondown{}Son independientes los sucesos ``gustar la salsa'' y ``gustar el merengue''? \textquestiondown{}Son compatibles?
\end{enumerate}
